\chapter{OPERATING MODES}
\label{ch:modes}

The software is arranged so that each sensor can be left out and the rest
still works. This chapter lists the modes that exist today, what each one
estimates, and how to start it.

\section{LOCALIZATION AND MAPPING MODES}

\begin{table}[htbp]
\centering
\caption{Localization and mapping modes}
\label{tab:modes}
\small
\begin{tabularx}{\linewidth}{>{\raggedright\arraybackslash}p{1.25in}>{\raggedright\arraybackslash}p{1.25in}>{\raggedright\arraybackslash}p{1.35in}X}
\toprule
Mode & Sensors & Start & Position estimate \\
\midrule
\textbf{A. gmapping SLAM} (default) & lidar, encoders, gyro &
  \texttt{robot.launch}; \texttt{start\_slam.sh} &
  \frm{map}: odometry corrected by scan matching; drift bounded while the
  map is consistent. \\
\textbf{B. hector SLAM with odometry TF} & lidar (odometry only as TF) &
  \texttt{start\_slam.sh hector} &
  \frm{map}: scan matching only; fails under fast turns. \\
\textbf{C. hector, lidar only} & lidar & \texttt{robot.launch}, then
  \texttt{start\_slam.sh lidar\_only} (or the A1 alone) &
  As B, with no wheel odometry at all; for carrying the robot. \\
\textbf{D. Dead reckoning} & encoders, gyro & \texttt{robot.launch lidar:=false} &
  \frm{odom} only: distance from encoders, heading from gyro with ZUPT;
  unbounded drift (\S\ref{sec:odom}). \\
\textbf{E. Encoder heading} & encoders & \texttt{use\_gyro\_heading: false} in
  \file{base.yaml} & \frm{odom} with heading from \eqref{eq:diffdrive};
  turns overestimated by $\chi$ (\S\ref{sec:skid}). Diagnostic only. \\
\textbf{F. GPS fusion} (parallel to A--D) & GPS, compass, gyro, encoders &
  \texttt{gps\_fusion.launch} on the laptop &
  \frm{gps\_map}: bounded by GPS when there is a fix; coasts on odometry
  with compass heading otherwise. \\
\textbf{G. Autonomous exploration} & lidar, encoders, gyro &
  \texttt{robot.launch}; \texttt{explore\_house.sh} on the laptop &
  As A; the robot drives itself to unexplored frontiers
  (\S\ref{sec:explore}). \\
\bottomrule
\end{tabularx}
\end{table}

Mode F runs alongside any of A--D because it publishes no transform. In
practice: indoors, run A and ignore F; outdoors, run F, and A as well if
there are walls or obstacles within the lidar's range.

\section{WHAT HAPPENS WHEN A SENSOR IS MISSING}

\begin{table}[htbp]
\centering
\caption{Degradation by sensor}
\label{tab:degrade}
\small
\begin{tabularx}{\linewidth}{>{\raggedright\arraybackslash}p{1.1in}X}
\toprule
Missing & Effect \\
\midrule
GPS fix & Normal indoors. \topic{/gps/fix} reports \texttt{NO\_FIX}; the EKF
  coasts. If the first fix comes later, navsat\_transform takes its datum
  then. \\
Compass & \topic{/imu/data} not published (\texttt{mag\_heading: false} or no
  field); the EKF has no absolute heading and navsat\_transform cannot set a
  datum, so mode F does not work. Odometry and SLAM are unaffected. \\
Gyroscope (IMU) & \texttt{bbai\_base} cannot start (the IMU is initialised
  unconditionally). Mode E would be the fallback but needs a code change to
  skip the IMU. \\
Encoders & Odometry reports no motion. Mode C still maps. ZUPT would treat the
  robot as parked whenever no command is active. \\
One encoder & Its side averages the remaining wheel with a stuck one, so that
  side reads half its travel; fix the channel list in \file{base.yaml}. \\
Lidar & Modes A--C unavailable; \texttt{robot.launch lidar:=false} gives D. \\
Motor 1 PRU firmware & \texttt{/dev/rpmsg\_pru32} does not appear (the
  base node logs a warning), \texttt{rc\_motor\_init} fails, and the motors
  cannot be driven. \\
Laptop & Everything on the board runs; no SLAM, no GPS fusion, no RViz. \\
Bluetooth controller & Motors receive no commands; the \SI{0.5}{s} timeout
  keeps them stopped. \\
\bottomrule
\end{tabularx}
\end{table}

\section{DRIVING MODES}

\begin{description}
\item[Manual (teleop).] Hold L1; left stick drives, right stick turns, R1
  adds turbo. Releasing L1 stops the robot.
\item[Chase, dry run.] \texttt{roslaunch bbai\_chase chase.launch}. Holding R2
  publishes only the preview topic; nothing moves.
\item[Chase, armed.] \texttt{chase.launch dry\_run:=false}. Holding R2 (with
  L1 released) lets the chase node drive toward the target. Needs the
  camera on its own power (\S\ref{sec:resettests}).
\item[Autonomous exploration.] \texttt{\textasciitilde/bbai\_slam/explore\_house.sh}
  on the laptop, with \texttt{robot.launch} on the board. The robot maps on
  its own (\S\ref{sec:explore}); Ctrl-C stops it and saves the map.
\item[Calibration.] \texttt{rosrun bbai\_base mag\_calibrate.py} drives
  \topic{/cmd\_vel} itself to spin two turns. \texttt{motor\_map\_test.py}
  drives each motor in turn and must be run with \texttt{robot.launch}
  stopped and the wheels lifted.
\end{description}

\section{AUTONOMOUS EXPLORATION}
\label{sec:explore}

Mode G lets the robot map the house without a driver. It runs entirely on
the laptop on top of gmapping (mode A) and uses three standard ROS
components:
\begin{description}
\item[gmapping] builds the occupancy grid and publishes \frm{map}$\to$\frm{odom}
  as in Chapter~\ref{ch:slam}.
\item[explore\_lite] finds \emph{frontiers}, the boundaries between known
  free cells and unknown cells \cite{yamauchi1997frontier}. Adjacent
  frontier cells are grouped, groups smaller than \SI{0.4}{m} are dropped,
  and each remaining frontier $f$ is scored by
  \begin{equation}\label{eq:frontier}
  c(f)=\lambda_d\,d(f)-\lambda_g\,\abs{f},
  \end{equation}
  with $d(f)$ the distance from the robot to the frontier, $\abs{f}$ its size,
  $\lambda_d=3$ (\texttt{potential\_scale}) and $\lambda_g=1$
  (\texttt{gain\_scale}). The cheapest frontier becomes the next goal, chosen
  again every \SI{3}{s}; a goal with no progress for \SI{30}{s} is
  blacklisted.
\item[move\_base] plans and drives to the goal. A global planner (navfn,
  Dijkstra's algorithm on the costmap, allowed through unknown space) gives a
  path, and the DWA local planner \cite{fox1997dynamic} follows it at
  \SI{5}{Hz}. DWA samples velocity pairs $(v,\omega)$ reachable within one
  control period under the acceleration limits, forward-simulates each for
  \SI{2}{s}, and picks
  \begin{equation}\label{eq:dwa}
  (v,\omega)^\ast=\argmin_{(v,\omega)}\
  \alpha\,d_{\mathrm{path}}+\beta\,d_{\mathrm{goal}}+\gamma\,c_{\mathrm{obs}},
  \end{equation}
  with $d_{\mathrm{path}}$ and $d_{\mathrm{goal}}$ the end point's distance
  from the global path and the goal, $c_{\mathrm{obs}}$ the highest costmap cost
  along the trajectory (trajectories through obstacles are discarded), and
  $\alpha=32$, $\beta=20$, $\gamma=0.05$.
\end{description}

\paragraph{Costmaps.} Obstacles come only from \topic{/scan}, marked up to
\SI{3}{m} and cleared by ray tracing to \SI{6}{m}. The footprint is
\SI{0.22}{m}$\times$\SI{0.20}{m} plus \SI{2}{cm} padding, and costs are
inflated out to \SI{0.25}{m} with
$c(r)=252\,e^{-8(r-r_i)}$ beyond the inscribed radius $r_i$. The local
costmap is a \SI{3}{m}$\times$\SI{3}{m} window in \frm{odom}; the global
costmap is the SLAM map.

\paragraph{Motion limits.} At most \SI{0.2}{m/s} forward, never backward
(the configuration avoids reversing because the lidar is the only obstacle
sensor), turn rates
\SIrange{0.6}{1.5}{rad/s}. The \SI{0.6}{rad/s} minimum exists because slower
turns stall the skid steer; with $\chi_c=\SI{0.46}{m}$ it is a duty of about
0.22. Each map is saved to \file{\textasciitilde/bbai\_slam/maps/} with a
timestamp when the script is stopped.

\paragraph{Limits and safety.}
\begin{itemize}
\item The RPLIDAR scans one horizontal plane about \SI{10}{cm} up. It cannot
  see stairs, drop-offs or anything lower than the scan plane (shoes, cables,
  pet bowls), and it may see table and chair legs only as thin points.
  \emph{Keep the robot away from stairs, or block them, while exploring.}
\item The camera stays off during exploration because of the power resets
  (\S\ref{sec:powerbudget}), so there is no visual obstacle or dog
  detection.
\item The deadman does not apply: move\_base publishes \topic{/cmd\_vel}
  itself at \SI{5}{Hz}. Holding L1 lets teleop interleave its commands, but
  the two will fight; Ctrl-C on the laptop is the stop, after which the base
  node's \SI{0.5}{s} timeout halts the motors.
\item Exploration inherits every SLAM weakness of Chapter~\ref{ch:slam}: if
  the map breaks, the frontiers are wrong and the robot drives toward them.
\end{itemize}
\influx{Exploration (configuration in \file{laptop/bbai\_slam/nav/}) has not
yet been run through a whole house.}
